% !TEX root = ../main.tex

\chapter{Deskripsi Persoalan}

\section{Latar Belakang}
\lipsum[1]

\section{Tujuan}
Tugas ini bertujuan untuk mengimplementasikan Feedforward Neural Network (FFNN) \textit{from scratch} tanpa menggunakan library deep learning yang sudah ada. Secara spesifik, tujuan tugas ini adalah:

\begin{enumerate}
    \item Memahami mekanisme \textit{forward propagation} pada FFNN.
    \item Memahami mekanisme \textit{backward propagation} dan perhitungan gradien menggunakan \textit{chain rule}.
    \item Menganalisis pengaruh berbagai hyperparameter terhadap performa model.
\end{enumerate}

\section{Dataset}
Dataset yang digunakan adalah \textbf{Global Student Placement \& Salary Dataset}, yang berisi data performa akademik mahasiswa beserta \textit{outcome} penempatan kerja dan gaji setelah lulus. Berikut deskripsi fitur-fitur yang digunakan:

\begin{table}[H]
    \centering
    \caption{Deskripsi Fitur Dataset}
    \begin{tabular}{lll}
        \toprule
        \textbf{Kolom} & \textbf{Deskripsi} & \textbf{Tipe} \\
        \midrule
        \texttt{cgpa} & Cumulative Grade Point Average (skala 4.0--10.0) & Float \\
        \texttt{backlogs} & Jumlah mata kuliah yang belum lulus & Integer \\
        \texttt{college\_tier} & Tingkatan institusi & Categorical \\
        \texttt{country} & Negara kelulusan & Categorical \\
        \texttt{university\_ranking\_band} & Kategori peringkat global universitas & Categorical \\
        \texttt{internship\_count} & Jumlah internship yang telah diselesaikan & Integer \\
        \texttt{aptitude\_score} & Skor tes kemampuan kuantitatif & Integer \\
        \texttt{communication\_score} & Skor kemampuan komunikasi & Integer \\
        \texttt{specialization} & Bidang akademik & Categorical \\
        \texttt{industry} & Industri tempat bekerja & Categorical \\
        \texttt{internship\_quality\_score} & Skor kualitas pengalaman internship & Float \\
        \texttt{placement\_status} & Status penempatan kerja (target) & Categorical \\
        \bottomrule
    \end{tabular}
    \label{tab:dataset}
\end{table}
